\section{Task Description}
Our research project aims to advance the Court Judgment Prediction and Explanation (CJPE) task, incorporating insights and methodologies from both \cite{malik-etal-2021-ildc} and \cite{vats-etal-2023-llms}. The CJPE task involves two key sub-tasks: Prediction and Explanation. These tasks are performed sequentially, addressing the critical need to predict legal judgments and provide explanations for these predictions. To provide a visual representation of our task framework, Figure \ref{fig:task-framework} illustrates the overall process of Court Judgment Prediction and Explanation (CJPE) as employed in our study. This figure includes the sequential steps of prediction and explanation. To demonstrate an Indian case structure in Table \ref{case-example} and how explanations are derived from judicial judgments, refer to Figure \ref{fig:annotated-example} in the Appendix. For context, you can view the original case text on which this annotation is based.\footnote{\url{https://indiankanoon.org/doc/97694707/}} This example showcases the detailed process by which our annotators have identified and extracted these critical parts, reflecting the essence of judicial reasoning in each case. Similarly, an Example of an Indian Case Structure Table \ref{case-example}.

\textbf{Prediction Task:} The core of the CJPE task is to predict the outcome of a legal case based on the case proceedings. Given a document \(D\) that includes the case proceedings from the Supreme Court of India (SCI), the task is to predict the decision \(y \in \{0, 1\}\), where `1' signifies the acceptance of the appeal or petition by the appellant or petitioner, and `0' indicates its rejection.

\textbf{Explanation Task:} The second part of the CJPE task involves explaining the predicted decision. Our approach is two-fold, integrating methodologies from both referenced papers:

1. \textbf{Identifying Key Sentences (ILDC for CJPE approach):} Similar to the approach in \cite{malik-etal-2021-ildc}, we focus on identifying and highlighting key sentences or segments within the case proceedings that significantly contributed to the predicted outcome. This method relies on extracting specific parts of the text that are directly related to the decision, providing an evidence-based explanation.

2. \textbf{Generating Abstract Reasoning (LLMs approach):} Drawing from the approach in \cite{vats-etal-2023-llms}, we attempt to generate more abstract reasoning for the prediction. This involves providing zero and few-shot examples to the LLMs to guide them in generating explanations that are not just tied to specific text excerpts but also encompass broader reasoning and legal principles.

Additionally, we introduce a novel aspect to this task by training the LLMs specifically for both prediction and explanation. This training is tailored to enable the models to understand and process legal texts more effectively, improving their capability to predict outcomes and generate relevant explanations.
